\section{Context Probes, Preference Samplers and Resource Details}
\label{app:implementation_details}

Table~\ref{tab:v23_budget} gives the model, training and per-run resource settings.

\subsection{Transition-derived context and local action-response probe}

The label-free context vector $z_t$ concatenates squashed reward exponential
averages $\tanh(\overline{\mathbf{r}}/3)$ (EMA rate $0.2$), a scalar local
action-response probe, and an eight-dimensional transition latent $e_\psi$
(online EMA rate $0.35$). Using zero-based state and action indices, the local
action-response probe records the short-step gain ratio $v = s'_1 / a_0$
whenever $|a_0| > 0.05$, $v$ is finite and $|v| < 10$, updating a causal EMA
$\bar{v}$ at rate $0.2$ from the first accepted sample and outputting
$\operatorname{clip}((\bar{v} - 1)/0.5, -3, 3)$ (zero prior to first acceptance).
This lightweight probe provides a fast responsiveness summary alongside reward
statistics and learned transition embeddings without assuming parametric system
identification.

The transition encoder maps $(s, a, \tanh(\mathbf{r}/3), s' - s)$ to an
eight-dimensional $\tanh$-bounded latent. Its auxiliary objective minimizes
mean-squared reconstruction error on $\tanh(\mathbf{r}/3)$ and scaled state
displacement $\tanh((s' - s)/3)$, adds an $\ell_2$ latent penalty with weight
$10^{-4}$, and applies a temporal contrastive loss (weight $0.2$, temperature
$0.2$, positive time window $256$ steps). Replay transitions retain their
collected context features during preference relabeling; no regime labels enter
training or evaluation.

\subsection{Preference samplers}

Four preference-sampling rules parameterize the portfolio across $m$ objectives
(Table~\ref{tab:operators}):
\begin{itemize}[leftmargin=1.5em,itemsep=2pt,topsep=3pt]
 \item \texttt{uniform}: draws $w \sim \mathrm{Dir}(\mathbf{1}_m)$ independently
       during both environment interaction and replay relabeling/candidate
       generation.
 \item \texttt{edge}: uses the $(m+1)$-point corners-plus-center set
       $\mathcal{E}_m = \{e_1, \ldots, e_m, \frac{1}{m}\mathbf{1}_m\}$, cycling
       deterministically by $(t-1) \bmod (m+1)$ during interaction and sampling
       uniformly from $\mathcal{E}_m$ in replay.
 \item \texttt{stratified}: for $m = 2$ (HalfCheetah), uses the 11-point uniform
       simplex grid $\{(i/10, 1 - i/10) : i = 0, \ldots, 10\}$ (cyclic in
       interaction, uniform in replay). For $m \geq 3$ (Hopper and Ant),
       \texttt{stratified} uses the $(m+1)$-point corner-center set
       $\mathcal{E}_m$, so slots C5 and R2 coincide functionally with C4 on
       three-objective tasks (yielding 13 functionally distinct settings on
       $m = 3$, with zero discrete mass on non-vertex edges of the
       three-objective simplex).
 \item \texttt{mixed}: during interaction at step $t$, selects \texttt{edge}
       when $t \equiv 0 \pmod 3$ ($1/3$), \texttt{stratified} on remaining even
       steps ($1/3$), and $\mathrm{Dir}(\mathbf{1}_m)$ on remaining odd steps
       ($1/3$). Each branch occurs twice per six consecutive steps.
       For three objectives, the first two branches together assign two
       thirds of interaction requests to the corners-plus-center set.
       In replay, each draw independently selects
       $\mathrm{Dir}(\mathbf{1}_m)$ or \texttt{stratified} with probability
       $0.5$ each.
\end{itemize}

\subsection{Block allocation hyperparameters and per-run resources}

During the first $5{,}000$ transitions, the initial coverage slot supplies
preferences while causal per-objective reward means and variances initialize
(EMA rate $1/1500$). Before each subsequent $1{,}500$-step block, per-objective
standard deviations are frozen with numerical floor $10^{-6}$, and completed
blocks update the active slot score via $q_o \leftarrow 0.87\,q_o + 0.13\,P_k$
prior to that step's gradient update (incomplete terminal blocks remain
unscored). After ten deterministic warmup blocks
(Appendix~\ref{app:operators}), the selector enforces minimum family quota
fractions $(0.36, 0.22, 0.18)$ by choosing the family with the largest positive
deficit; otherwise it samples a family via a softmax over the mean of its two
highest slot scores (temperature $0.22$, uniform exploration mixture $0.16$)
and then samples a slot within that family (temperature $0.28$, exploration
mixture $0.16$). Within-family sampling still uses slot scores during family
deficit correction. The score-neutral control passes zero scores to routing
and keeps the quotas, warmup blocks and exploration mixtures.
At 75k interactions there are 46 complete feedback blocks: ten scheduled
warmup blocks and 36 subsequent blocks. The last 1,000 training steps use
a selected slot but provide no complete-block score.

\begin{table}[htbp]
\centering\small
\caption{Hyperparameters and per-run computational settings for the 75k
sample-constrained benchmark. Hidden widths apply to two-layer MLPs; the 50k
archive study uses identical network and optimizer settings while retaining 19
periodic actor snapshots.}
\label{tab:v23_budget}
\begin{tabular}{lcc}
\toprule
Quantity & BAMOR & CAPQL (Native / Full-Simplex)\\
\midrule
Training interactions & 75,000 & 75,000\\
Learning starts / replay minibatch & 5,000 / 256 & 5,000 / 256\\
Actor and critic hidden width & 64 & 256\\
Replay buffer capacity & 500,000 & 500,000\\
Learning rate / target Polyak rate & $3\times 10^{-4}$ / $0.005$ & $3\times 10^{-4}$ / $0.005$\\
Evaluated actors / preference grid $|W|$ & 1 / 21 & 1 / 21\\
Evaluation episodes per weight and mode & 3 & 3\\
Preference copies per update ($C_o$) & 1--3 & 1\\
Critic / actor candidates ($K_o^Q$ / $K_o^\pi$) & 1--5 / 1--5 & 1 / 1\\
Training preference distribution & 18-slot portfolio & $22.5^\circ$ cone / $\mathrm{Dir}(\mathbf{1})$\\
\bottomrule
\end{tabular}
\end{table}
